\section{Discussion and Validation Implications}
\label{sec:discussion}
\subsection{What the Benchmark Establishes}
The fixture makes a useful logical distinction: network loss and declared workload loss need not coincide. The unsupported wearable request provides a counterexample to universal local continuity within the corpus. The degraded-Wi-Fi snapshots provide a second counterexample, showing that a stress name alone cannot demonstrate degraded behavior. Together, these findings support a benchmark-design lesson: report the exercised states and unsupported workload combinations, as well as successful labels.

The seeded experiment adds reproducible variability and explicit completion and occupancy outcomes. It does not validate the classifier, because it samples a separate minimum-useful Boolean from class priors. Its outcomes can be analytically understood from those priors. Increasing the seed count reduces Monte Carlo uncertainty about the chosen model; it does not reduce uncertainty about whether that model represents a real device or workload. The relatively narrow intervals should be interpreted in that light.

The endpoint audit also changes the interpretation of broad metric panels. Reporting reliability, loss, outage, and QoE together can suggest several confirming observations even when the code computes each from the same Boolean sequence. An audit should trace each metric to its generating observations and retain distinct measurements only where their relationship is explicit. This study exposes the redundancy rather than assigning a composite score or selecting the most favorable endpoint.

\subsection{Missing Comparisons and Prospective Study Design}
The current repository does not execute terrestrial-only, cloud-only, local-only, or oracle continuity controllers in this experiment. Literature comparisons in Section~\ref{sec:related} are conceptual. A fair prospective study would give each executed controller the same workload arrivals, sensor information, device budgets, outage traces, switching costs, and recovery opportunities. If an oracle receives future information, it should be labeled an unattainable reference bound and analyzed separately. These conditions are future design requirements, not guarantees established by the current code.

A physical study should define a minimum useful outcome per workload before collecting comparison data. For a cached lesson this may include readable content and saved progress; for coding it may include correct local execution of a named program; for payment it must address authorization and durable transaction semantics, so a generic offline-capability flag is insufficient. Numerical latency, deadline, accuracy, memory, energy, and data-freshness thresholds should be justified through task requirements or a separate user study. No such physical threshold study has been completed here.

A second prospective stage should run those workloads during controlled bearer loss, congestion, device resource pressure, and restoration, recording timestamps and task outcomes. It should compare classification against observed utility, including connected failures and false declarations of offline support. Calibration data and held-out validation traces must be separated. Bursty outages and correlated demand should replace independent stress only under a new versioned protocol. This would connect the current semantic corpus to a measured twin-state construction without retroactively changing the frozen benchmark.

\subsection{Relation to the Three-Paper Research Program}
This article provides workload semantics and an auditable reporting boundary. A separate cross-layer-control study asks how radio and orchestration decisions change modeled performance; a separate terrestrial/non-terrestrial study asks when fallback paths are beneficial or costly. Those decision problems require their own implemented controllers and experiments. Shared motivation or vocabulary does not prove independent scholarly novelty. Reuse of this corpus in either companion study must cite its exact version and must not count the same outcomes again as independent evidence.

The earlier physical-measurement roadmap and the later minimum-useful-service roadmap remain compatible as sequential research stages: this benchmark defines what should be observed, while field campaigns could determine whether its assumptions hold. Today the evidence supports the benchmark stage. It does not establish physical channel characterization, synchronized twin-state fidelity, operational 6G performance, or any university-affiliated research activity.
